% Ad Atticum 13.33 (ad-atticum-13-33)
% Source: https://raw.githubusercontent.com/PerseusDL/canonical-latinLit/master/data/phi0474/phi057/phi0474.phi057.perseus-lat1.xml
% Fetched by scripts/fetch_latin.py

% Dateline (Perseus): Scr. in Tusculano iii Non. Iun. a. 709 (45).

\ciceroLetterOpener{CICERO ATTICO salutem}

\ciceroSection{1} O neglegentiam miram! semelne putas mihi dixisse Balbum et Faberium professionem relatam? qui etiam eorum iussu miserim qui profiteretur. ita enim oportere dicebant. professus est Philotimus libertus. nosti , credo, librarium.

\ciceroSection{2} sed scribes et quidem confectum. ad Faberium, ut tibi placet, litteras misi, cum Balbo autem puto te aliquid fecisse †H in Capitolio†. in Vergilio mihi nulla est δυσωπία . nec enim eius causa sane debeo et, si emero, quid erit quod postulet? sed videbis ne is tum sit in Africa ut Caelius. de nomine tu videbis cum Cispio; sed si Plancus destinat, tum habet res difficultatem. te ad me venire uterque nostrum cupit; sed ista res nullo modo relinquenda est. Othonem quod speras posse vinci sane bene narras. de aestimatione, ut scribis, cum agere coeperimus; etsi nihil scripsit nisi de modo agri. cum Pisone, si quid poterit. Dicaearchi librum accepi et καταβάσεωσ †exspecto†.

\ciceroSection{3} * * * negotium dederis, reperiet ex eo libro in quo sunt servatus consulta Cn. Cornelio L. Mummio coss. de Tuditano autem quod putas, εὔλογον est tum illum, quoniam fuit ad Corinthum (non enim temere dixit Hortensius), aut quaestorem aut tribunum mil. fuisse, idque potius credo. tu de Antiocho scire poteris †vide etiam†, quo anno quaestor aut tribunus mil. fuerit; si neutrum †ea de† in praefectis an in contubernalibus fuerit, modo fuerit in eo bello.
